\section{Мета роботи та завдання}

\textbf{Мета роботи:}
\begin{enumerate}
    \item Встановити необхідне програмне забезпечення для програмування на ЯП \texttt{Java};
    \item Ознайомитися з базовим синтаксисом мови та основними операторами;
    \item Вивчити приклади угоди щодо коду для ЯП \texttt{Java}.
\end{enumerate}

\skipline
\textbf{Завдання на лабораторну роботу:}
\begin{enumerate}
    \item 
    Напишіть метод, який приймає на вхід масив цілих чисел (довжиною 2 і більше) і 
    повертає true, якщо в масиві кожен елемент більше або дорівнює попередньому. Інакше він повертає false.
    \item 
    Напишіть метод, 
    який виводить на екран числа від 1 до 100. При цьому замість чисел, кратних трьох, 
    програма повинна виводити слово Fizz, а замість чисел, кратних п'яти - слово Buzz. 
    Якщо число кратне і 3, і 5, програма повинна виводити слово «FizzBuzz».
    \item 
    Реалізуйте метод, який приймає на вхід масив цілих чисел (довжиною 2 або більше) і повертає true, 
    якщо масив можна розділити так, щоб сума чисел в обох частинах була рівною і false інакше.
\end{enumerate}
